% Options for packages loaded elsewhere
% Options for packages loaded elsewhere
\PassOptionsToPackage{unicode}{hyperref}
\PassOptionsToPackage{hyphens}{url}
%
\documentclass[
  14pt,
  ignorenonframetext,
  aspectratio=169,
]{beamer}
\newif\ifbibliography
\usepackage{pgfpages}
\setbeamertemplate{caption}[numbered]
\setbeamertemplate{caption label separator}{: }
\setbeamercolor{caption name}{fg=normal text.fg}
\beamertemplatenavigationsymbolsempty
% remove section numbering
\setbeamertemplate{part page}{
  \centering
  \begin{beamercolorbox}[sep=16pt,center]{part title}
    \usebeamerfont{part title}\insertpart\par
  \end{beamercolorbox}
}
\setbeamertemplate{section page}{
  \centering
  \begin{beamercolorbox}[sep=12pt,center]{section title}
    \usebeamerfont{section title}\insertsection\par
  \end{beamercolorbox}
}
\setbeamertemplate{subsection page}{
  \centering
  \begin{beamercolorbox}[sep=8pt,center]{subsection title}
    \usebeamerfont{subsection title}\insertsubsection\par
  \end{beamercolorbox}
}
% Prevent slide breaks in the middle of a paragraph
\widowpenalties 1 10000
\raggedbottom
\usepackage{iftex}
\ifPDFTeX
  \usepackage[T1]{fontenc}
  \usepackage[utf8]{inputenc}
  \usepackage{textcomp} % provide euro and other symbols
\else % if luatex or xetex
  \usepackage{unicode-math} % this also loads fontspec
  \defaultfontfeatures{Scale=MatchLowercase}
  \defaultfontfeatures[\rmfamily]{Ligatures=TeX,Scale=1}
\fi
\usepackage{lmodern}

\usetheme[]{monash}
\ifPDFTeX\else
  % xetex/luatex font selection
\fi
% Use upquote if available, for straight quotes in verbatim environments
\IfFileExists{upquote.sty}{\usepackage{upquote}}{}
\IfFileExists{microtype.sty}{% use microtype if available
  \usepackage[]{microtype}
  \UseMicrotypeSet[protrusion]{basicmath} % disable protrusion for tt fonts
}{}
\makeatletter
\@ifundefined{KOMAClassName}{% if non-KOMA class
  \IfFileExists{parskip.sty}{%
    \usepackage{parskip}
  }{% else
    \setlength{\parindent}{0pt}
    \setlength{\parskip}{6pt plus 2pt minus 1pt}}
}{% if KOMA class
  \KOMAoptions{parskip=half}}
\makeatother

\usepackage{color}
\usepackage{fancyvrb}
\newcommand{\VerbBar}{|}
\newcommand{\VERB}{\Verb[commandchars=\\\{\}]}
\DefineVerbatimEnvironment{Highlighting}{Verbatim}{commandchars=\\\{\}}
% Add ',fontsize=\small' for more characters per line
\usepackage{framed}
\definecolor{shadecolor}{RGB}{248,248,248}
\newenvironment{Shaded}{\begin{snugshade}}{\end{snugshade}}
\newcommand{\AlertTok}[1]{\textcolor[rgb]{0.94,0.16,0.16}{#1}}
\newcommand{\AnnotationTok}[1]{\textcolor[rgb]{0.56,0.35,0.01}{\textbf{\textit{#1}}}}
\newcommand{\AttributeTok}[1]{\textcolor[rgb]{0.13,0.29,0.53}{#1}}
\newcommand{\BaseNTok}[1]{\textcolor[rgb]{0.00,0.00,0.81}{#1}}
\newcommand{\BuiltInTok}[1]{#1}
\newcommand{\CharTok}[1]{\textcolor[rgb]{0.31,0.60,0.02}{#1}}
\newcommand{\CommentTok}[1]{\textcolor[rgb]{0.56,0.35,0.01}{\textit{#1}}}
\newcommand{\CommentVarTok}[1]{\textcolor[rgb]{0.56,0.35,0.01}{\textbf{\textit{#1}}}}
\newcommand{\ConstantTok}[1]{\textcolor[rgb]{0.56,0.35,0.01}{#1}}
\newcommand{\ControlFlowTok}[1]{\textcolor[rgb]{0.13,0.29,0.53}{\textbf{#1}}}
\newcommand{\DataTypeTok}[1]{\textcolor[rgb]{0.13,0.29,0.53}{#1}}
\newcommand{\DecValTok}[1]{\textcolor[rgb]{0.00,0.00,0.81}{#1}}
\newcommand{\DocumentationTok}[1]{\textcolor[rgb]{0.56,0.35,0.01}{\textbf{\textit{#1}}}}
\newcommand{\ErrorTok}[1]{\textcolor[rgb]{0.64,0.00,0.00}{\textbf{#1}}}
\newcommand{\ExtensionTok}[1]{#1}
\newcommand{\FloatTok}[1]{\textcolor[rgb]{0.00,0.00,0.81}{#1}}
\newcommand{\FunctionTok}[1]{\textcolor[rgb]{0.13,0.29,0.53}{\textbf{#1}}}
\newcommand{\ImportTok}[1]{#1}
\newcommand{\InformationTok}[1]{\textcolor[rgb]{0.56,0.35,0.01}{\textbf{\textit{#1}}}}
\newcommand{\KeywordTok}[1]{\textcolor[rgb]{0.13,0.29,0.53}{\textbf{#1}}}
\newcommand{\NormalTok}[1]{#1}
\newcommand{\OperatorTok}[1]{\textcolor[rgb]{0.81,0.36,0.00}{\textbf{#1}}}
\newcommand{\OtherTok}[1]{\textcolor[rgb]{0.56,0.35,0.01}{#1}}
\newcommand{\PreprocessorTok}[1]{\textcolor[rgb]{0.56,0.35,0.01}{\textit{#1}}}
\newcommand{\RegionMarkerTok}[1]{#1}
\newcommand{\SpecialCharTok}[1]{\textcolor[rgb]{0.81,0.36,0.00}{\textbf{#1}}}
\newcommand{\SpecialStringTok}[1]{\textcolor[rgb]{0.31,0.60,0.02}{#1}}
\newcommand{\StringTok}[1]{\textcolor[rgb]{0.31,0.60,0.02}{#1}}
\newcommand{\VariableTok}[1]{\textcolor[rgb]{0.00,0.00,0.00}{#1}}
\newcommand{\VerbatimStringTok}[1]{\textcolor[rgb]{0.31,0.60,0.02}{#1}}
\newcommand{\WarningTok}[1]{\textcolor[rgb]{0.56,0.35,0.01}{\textbf{\textit{#1}}}}

\usepackage{longtable,booktabs,array}
\newcounter{none} % for unnumbered tables
\usepackage{calc} % for calculating minipage widths
\usepackage{caption}
% Make caption package work with longtable
\makeatletter
\def\fnum@table{\tablename~\thetable}
\makeatother
\usepackage{graphicx}
\makeatletter
\newsavebox\pandoc@box
\newcommand*\pandocbounded[1]{% scales image to fit in text height/width
  \sbox\pandoc@box{#1}%
  \Gscale@div\@tempa{\textheight}{\dimexpr\ht\pandoc@box+\dp\pandoc@box\relax}%
  \Gscale@div\@tempb{\linewidth}{\wd\pandoc@box}%
  \ifdim\@tempb\p@<\@tempa\p@\let\@tempa\@tempb\fi% select the smaller of both
  \ifdim\@tempa\p@<\p@\scalebox{\@tempa}{\usebox\pandoc@box}%
  \else\usebox{\pandoc@box}%
  \fi%
}
% Set default figure placement to htbp
\def\fps@figure{htbp}
\makeatother





\setlength{\emergencystretch}{3em} % prevent overfull lines

\providecommand{\tightlist}{%
  \setlength{\itemsep}{0pt}\setlength{\parskip}{0pt}}



 


% Colors
\definecolor{shadecolor}{RGB}{225,225,225}
\setbeamercolor{description item}{fg=Orange}
\definecolor{MonashBlue}{RGB}{66,109,152}
\definecolor{burntorange}{rgb}{0.8, 0.33, 0.0}

% Packages
\usepackage{amsmath, bm, amssymb, amsthm, mathrsfs,pifont}
\usepackage{url}
\usepackage{multirow, booktabs, float, textcmds, siunitx}
\usepackage{bm,booktabs,animate,ragged2e,multicol,microtype,hyperref}
\usepackage{fontawesome5}

% Figures
\graphicspath{{figs/}}

% Fonts
\fontsize{13}{15}\sf
\setbeamerfont{title}{series=\bfseries,parent=structure,size=\fontsize{24}{24}}
\ifcsname Shaded\endcsname
  \definecolor{shadecolor}{RGB}{225,225,225}
  \renewenvironment{Shaded}{\vspace*{0.15cm}\color{black}\fontsize{10}{10}\sf\begin{snugshade}\color{black}}{\end{snugshade}}
\fi

% gt dependencies
\usepackage{longtable,caption,setspace}
\captionsetup{font={small,stretch=0.80}}

% My definitions

\def\E{\text{E}}
\def\V{\text{Var}}
\def\up#1{\raisebox{-0.3cm}{#1}}
\def\pred#1#2#3{\hat{#1}_{#2|#3}}
\def\damped{$_\text{d}$}
\def\h+{h_{m}^{+}}
\def\str#1{\rlap{#1}\textcolor{red}{\rule{1cm}{0.1cm}}}

\def\st#1{\rlap{#1}\textcolor{red}{\rule{1cm}{0.1cm}}}
\def\bY{\bm{y}}
\def\by{\bm{y}}
\def\bS{\bm{S}}
\def\bI{\text{\rm\textbf{I}}}
\def\bbeta{\bm{\beta}}
\def\bSigma{\bm{\Sigma}}
\def\bW{\bm{\Sigma}}
\def\Var{\text{Var}}
\def\var{\text{Var}}
\def\bOmega{\bm{\Omega}}
\def\bLambda{\bm{\Lambda}}
\let\mc\multicolumn
\def\hl{\color[RGB]{230, 172, 0}}

\def\forecast{\begin{alertblock}{}\fontsize{10}{11}\sf
 A forecast is an estimate of the probability distribution of a variable to be observed in the future.
\end{alertblock}}

\def\simfutures{\begin{textblock}{2.9}(12.8,7.7)
\begin{block}{}\fontsize{10}{11}\sf
Simulated futures from an ETS model
\end{block}\end{textblock}}

\setbeamertemplate{title page}
{\placefig{0}{-0.5}{width=1.01\paperwidth,height=1.45\paperheight}{africa1.jpg}
\begin{textblock}{8.5}(.5,0.2)
  \raggedright\fontsize{20}{20}\selectfont\bfseries\sffamily\textcolor{white}{\inserttitle}\\[0.2cm]
  \fontsize{12}{12}\selectfont\bfseries\sffamily\textcolor{white}{\insertsubtitle}
  \end{textblock}
\placefig{.5}{6.2}{width=2.4cm}{tsibble.png}
\placefig{2.9}{6.2}{width=2.4cm}{feasts.png}
\placefig{5.3}{6.2}{width=2.4cm}{fable.png}
\begin{textblock}{7.5}(10.6,-.1)
  {\fontsize{9}{9}\sf\color{white}https://workshop.f4sg.org/africast/}
\end{textblock}
}

% Tikz plots
\usepackage{tikz}
\usetikzlibrary{trees,shapes,arrows,matrix,shadows,positioning,calc}
\tikzstyle{decision} = [diamond, draw, fill=blue!20,
    text width=4.5em, text badly centered, node distance=4cm, inner sep=0pt]
\tikzstyle{block} = [rectangle, draw, fill=blue!20,
    text width=5cm, text centered, rounded corners, minimum height=4em]
\tikzstyle{line} = [draw, thick, -latex']
\tikzstyle{cloud} = [draw, ellipse,fill=red!20, node distance=3cm,
    minimum height=2em, text centered]
\tikzstyle{connector} = [->,thick]

\tikzset{
  basic/.style  = {draw, text width=2cm, font=\sffamily, rectangle},
  root/.style   = {basic, text width=3cm, rounded corners=2pt, thin, align=center, fill=red!30},
  level 2a/.style = {basic, rounded corners=2pt, thin,align=center, fill=blue!50, text width=7em},
  level 2b/.style = {basic, rounded corners=2pt, thin,align=center, fill=green!50, text width=7em},
  level 3a/.style = {basic, rounded corners=2pt, thin, align=center, fill=blue!30, text width=4em},
  level 3b/.style = {basic, rounded corners=2pt, thin, align=center, fill=green!30, text width=4em},
  level 4a/.style = {basic, rounded corners=2pt, thin, align=left, fill=blue!10, text width=3.5em},
  level 4b/.style = {basic, rounded corners=2pt, thin, align=left, fill=green!10, text width=3.5em}
}

% % BIBLIOGRAPHIES
% \usepackage[style=authoryear,bibencoding=utf8,minnames=1,maxnames=4, maxbibnames=99,natbib=true,dashed=false,terseinits=true,giveninits=true,uniquename=false,uniquelist=false,labeldate=true,doi=false, isbn=false, natbib=true,backend=biber]{biblatex}

% \DeclareFieldFormat{url}{\url{#1}}
% \DeclareFieldFormat[article]{pages}{#1}
% \DeclareFieldFormat[inproceedings]{pages}{\lowercase{pp.}#1}
% \DeclareFieldFormat[incollection]{pages}{\lowercase{pp.}#1}
% \DeclareFieldFormat[article]{volume}{\mkbibbold{#1}}
% \DeclareFieldFormat[article]{number}{\mkbibparens{#1}}
% \DeclareFieldFormat[article]{title}{\MakeCapital{#1}}
% \DeclareFieldFormat[article]{url}{}
% \DeclareFieldFormat[Techreport]{Url}{}
% \DeclareFieldFormat[book]{url}{}
% \DeclareFieldFormat[inbook]{url}{}
% \DeclareFieldFormat[incollection]{url}{}
% \DeclareFieldFormat[inproceedings]{url}{}
% \DeclareFieldFormat[inproceedings]{title}{#1}
% \DeclareFieldFormat{shorthandwidth}{#1}
% %\DeclareFieldFormat{extrayear}{}
% % No dot before number of articles
% \usepackage{xpatch}
% \xpatchbibmacro{volume+number+eid}{\setunit*{\adddot}}{}{}{}
% % Remove In: for an article.
% \renewbibmacro{in:}{%
%   \ifentrytype{article}{}{%
%   \printtext{\bibstring{in}\intitlepunct}}}

% \AtEveryBibitem{\clearfield{month}}
% \AtEveryCitekey{\clearfield{month}}
% \AtBeginBibliography{\fontsize{11}{11}\sf}
\setbeamertemplate{frametitle continuation}{}
\makeatletter
\@ifpackageloaded{caption}{}{\usepackage{caption}}
\AtBeginDocument{%
\ifdefined\contentsname
  \renewcommand*\contentsname{Table of contents}
\else
  \newcommand\contentsname{Table of contents}
\fi
\ifdefined\listfigurename
  \renewcommand*\listfigurename{List of Figures}
\else
  \newcommand\listfigurename{List of Figures}
\fi
\ifdefined\listtablename
  \renewcommand*\listtablename{List of Tables}
\else
  \newcommand\listtablename{List of Tables}
\fi
\ifdefined\figurename
  \renewcommand*\figurename{Figure}
\else
  \newcommand\figurename{Figure}
\fi
\ifdefined\tablename
  \renewcommand*\tablename{Table}
\else
  \newcommand\tablename{Table}
\fi
}
\@ifpackageloaded{float}{}{\usepackage{float}}
\floatstyle{ruled}
\@ifundefined{c@chapter}{\newfloat{codelisting}{h}{lop}}{\newfloat{codelisting}{h}{lop}[chapter]}
\floatname{codelisting}{Listing}
\newcommand*\listoflistings{\listof{codelisting}{List of Listings}}
\makeatother
\makeatletter
\makeatother
\makeatletter
\@ifpackageloaded{caption}{}{\usepackage{caption}}
\@ifpackageloaded{subcaption}{}{\usepackage{subcaption}}
\makeatother

\usepackage{bookmark}
\IfFileExists{xurl.sty}{\usepackage{xurl}}{} % add URL line breaks if available
\urlstyle{same}
% fallback for those not using the hyperref driver hyperxmp:
\makeatletter
\@ifundefined{xmpquote}{\newcommand{\xmpquote}[1]{#1}}{}
\makeatother
\hypersetup{
  pdftitle={Africast-Time Series Analysis \& Forecasting Using R},
  hidelinks,
  pdfcreator={LaTeX via pandoc}}


\title{\texorpdfstring{Africast-Time Series Analysis \& Forecasting
Using R}{Africast-Time Series Analysis \& Forecasting Using R}}
\subtitle{\texorpdfstring{2. Time series patterns and basic
graphics}{2. Time series patterns and basic graphics}}
\author{}
\date{7 October 2026}

\begin{document}
\frame{\titlepage}


\begin{frame}{Outline}
\protect\phantomsection\label{outline}
\vspace*{0.7cm}\tableofcontents
\end{frame}

\section{Time series Patterns}\label{time-series-patterns}

\begin{frame}{Key patterns of time series}
\protect\phantomsection\label{key-patterns-of-time-series}
\begin{itemize}
\item
  Level
\item
  Underlying trend
\item
  Seasonal/cycle
\item
  Autocorrelation
\item
  Unpredictable patterns/Noise
\item
  Different types of events and driving factors (i.e.~predictors) may
  affect the time series
\end{itemize}
\end{frame}

\begin{frame}{Time series patterns}
\protect\phantomsection\label{time-series-patterns-1}
\begin{description}
\tightlist
\item[Level]
The \emph{level} of a time series describes the centre of the series.
\item[Trend]
A \emph{trend} describes predictable increases or decreases in the level
of a series.
\item[Seasonal]
\emph{Seasonality} is a consistent pattern that repeats over a fixed
period of time. pattern exists when a series is influenced by seasonal
factors (e.g., the quarter of the year, the month, or day of the week).
\item[Cyclic]
pattern exists when data exhibit rises and falls that are
\emph{not of fixed period} (duration usually of at least 2 years).
\end{description}
\end{frame}

\begin{frame}{Level}
\protect\phantomsection\label{level}
\begin{center}
\includegraphics[width=1\linewidth,height=\textheight,keepaspectratio]{01_graphics_files/figure-beamer/level-1.pdf}
\end{center}
\end{frame}

\begin{frame}{Trend}
\protect\phantomsection\label{trend}
\begin{center}
\includegraphics[width=1\linewidth,height=\textheight,keepaspectratio]{01_graphics_files/figure-beamer/trend-1.pdf}
\end{center}
\end{frame}

\begin{frame}{Seasonality}
\protect\phantomsection\label{seasonality}
\begin{center}
\includegraphics[width=1\linewidth,height=\textheight,keepaspectratio]{01_graphics_files/figure-beamer/additive-seasonality-1.pdf}
\end{center}
\end{frame}

\begin{frame}{Additive versus multiplicative seasonality}
\protect\phantomsection\label{additive-versus-multiplicative-seasonality}
\begin{center}
\includegraphics[width=1\linewidth,height=\textheight,keepaspectratio]{01_graphics_files/figure-beamer/trend-both-seasonality-1.pdf}
\end{center}
\end{frame}

\begin{frame}{Cycles}
\protect\phantomsection\label{cycles}
\begin{center}
\includegraphics[width=1\linewidth,height=\textheight,keepaspectratio]{01_graphics_files/figure-beamer/cement-1.pdf}
\end{center}
\end{frame}

\section{Time plots}\label{time-plots}

\begin{frame}[fragile]{Time plots}
\protect\phantomsection\label{time-plots-1}
\small

\begin{Shaded}
\begin{Highlighting}[]
\NormalTok{ansett }\SpecialCharTok{\%\textgreater{}\%}
  \FunctionTok{filter}\NormalTok{(Airports}\SpecialCharTok{==}\StringTok{"MEL{-}SYD"}\NormalTok{, Class}\SpecialCharTok{==}\StringTok{"Economy"}\NormalTok{) }\SpecialCharTok{\%\textgreater{}\%}
  \FunctionTok{autoplot}\NormalTok{(Passengers)}
\end{Highlighting}
\end{Shaded}

\pandocbounded{\includegraphics[keepaspectratio]{01_graphics_files/figure-beamer/unnamed-chunk-1-1.pdf}}
\end{frame}

\begin{frame}[fragile]{Time plots}
\protect\phantomsection\label{time-plots-2}
\small

\begin{Shaded}
\begin{Highlighting}[]
\NormalTok{PBS }\SpecialCharTok{\%\textgreater{}\%} \FunctionTok{filter}\NormalTok{(ATC2 }\SpecialCharTok{==} \StringTok{"A10"}\NormalTok{) }\SpecialCharTok{\%\textgreater{}\%}
  \FunctionTok{summarise}\NormalTok{(}\AttributeTok{Cost =} \FunctionTok{sum}\NormalTok{(Cost)}\SpecialCharTok{/}\FloatTok{1e6}\NormalTok{) }\SpecialCharTok{\%\textgreater{}\%} \FunctionTok{autoplot}\NormalTok{(Cost) }\SpecialCharTok{+}
  \FunctionTok{ylab}\NormalTok{(}\StringTok{"$ million"}\NormalTok{) }\SpecialCharTok{+} \FunctionTok{xlab}\NormalTok{(}\StringTok{"Year"}\NormalTok{) }\SpecialCharTok{+}
  \FunctionTok{ggtitle}\NormalTok{(}\StringTok{"Antidiabetic drug sales"}\NormalTok{)}
\end{Highlighting}
\end{Shaded}

\pandocbounded{\includegraphics[keepaspectratio]{01_graphics_files/figure-beamer/a10-plot-1.pdf}}
\end{frame}

\section{Seasonal plots}\label{seasonal-plots}

\begin{frame}[fragile]{Seasonal plots}
\protect\phantomsection\label{seasonal-plots-1}
\begin{itemize}
\tightlist
\item
  Data plotted against the individual ``seasons'' in which the data were
  observed. (In this case a ``season'' is a month.)
\item
  Something like a time plot except that the data from each season are
  overlapped.
\item
  Enables the underlying seasonal pattern to be seen more clearly, and
  also allows any substantial departures from the seasonal pattern to be
  easily identified.
\item
  In R: \texttt{gg\_season()} (ggtime)
\end{itemize}
\end{frame}

\begin{frame}[fragile]{Quarterly Australian Beer Production}
\protect\phantomsection\label{quarterly-australian-beer-production}
\begin{Shaded}
\begin{Highlighting}[]
\NormalTok{beer }\OtherTok{\textless{}{-}}\NormalTok{ aus\_production }\SpecialCharTok{|\textgreater{}}
  \FunctionTok{select}\NormalTok{(Quarter, Beer) }\SpecialCharTok{|\textgreater{}}
  \FunctionTok{filter}\NormalTok{(}\FunctionTok{year}\NormalTok{(Quarter) }\SpecialCharTok{\textgreater{}=} \FunctionTok{year}\NormalTok{(}\DecValTok{1992}\NormalTok{L))}
\NormalTok{beer }\SpecialCharTok{|\textgreater{}} \FunctionTok{autoplot}\NormalTok{(Beer)}
\end{Highlighting}
\end{Shaded}

\pandocbounded{\includegraphics[keepaspectratio]{01_graphics_files/figure-beamer/unnamed-chunk-2-1.pdf}}
\end{frame}

\begin{frame}[fragile]{Quarterly Australian Beer Production}
\protect\phantomsection\label{quarterly-australian-beer-production-1}
\begin{Shaded}
\begin{Highlighting}[]
\NormalTok{beer }\SpecialCharTok{|\textgreater{}} \FunctionTok{gg\_season}\NormalTok{(Beer, }\AttributeTok{labels =} \StringTok{"right"}\NormalTok{)}
\end{Highlighting}
\end{Shaded}

\pandocbounded{\includegraphics[keepaspectratio]{01_graphics_files/figure-beamer/unnamed-chunk-3-1.pdf}}
\end{frame}

\begin{frame}[fragile]{Multiple seasonal periods}
\protect\phantomsection\label{multiple-seasonal-periods}
\fontsize{10}{11}\sf

\begin{Shaded}
\begin{Highlighting}[]
\NormalTok{vic\_elec}
\end{Highlighting}
\end{Shaded}

\begin{verbatim}
# A tsibble: 52,608 x 5 [30m] <Australia/Melbourne>
   Time                Demand Temperature Date       Holiday
   <dttm>               <dbl>       <dbl> <date>     <lgl>  
 1 2012-01-01 00:00:00  4383.        21.4 2012-01-01 TRUE   
 2 2012-01-01 00:30:00  4263.        21.0 2012-01-01 TRUE   
 3 2012-01-01 01:00:00  4049.        20.7 2012-01-01 TRUE   
 4 2012-01-01 01:30:00  3878.        20.6 2012-01-01 TRUE   
 5 2012-01-01 02:00:00  4036.        20.4 2012-01-01 TRUE   
 6 2012-01-01 02:30:00  3866.        20.2 2012-01-01 TRUE   
 7 2012-01-01 03:00:00  3694.        20.1 2012-01-01 TRUE   
 8 2012-01-01 03:30:00  3562.        19.6 2012-01-01 TRUE   
 9 2012-01-01 04:00:00  3433.        19.1 2012-01-01 TRUE   
10 2012-01-01 04:30:00  3359.        19.0 2012-01-01 TRUE   
# i 52,598 more rows
\end{verbatim}
\end{frame}

\begin{frame}[fragile]{Multiple seasonal periods}
\protect\phantomsection\label{multiple-seasonal-periods-1}
\begin{Shaded}
\begin{Highlighting}[]
\NormalTok{vic\_elec }\SpecialCharTok{|\textgreater{}} \FunctionTok{gg\_season}\NormalTok{(Demand)}
\end{Highlighting}
\end{Shaded}

\pandocbounded{\includegraphics[keepaspectratio]{01_graphics_files/figure-beamer/unnamed-chunk-5-1.png}}
\end{frame}

\begin{frame}[fragile]{Multiple seasonal periods}
\protect\phantomsection\label{multiple-seasonal-periods-2}
\begin{Shaded}
\begin{Highlighting}[]
\NormalTok{vic\_elec }\SpecialCharTok{|\textgreater{}} \FunctionTok{gg\_season}\NormalTok{(Demand, }\AttributeTok{period =} \StringTok{"week"}\NormalTok{)}
\end{Highlighting}
\end{Shaded}

\pandocbounded{\includegraphics[keepaspectratio]{01_graphics_files/figure-beamer/unnamed-chunk-6-1.png}}
\end{frame}

\begin{frame}[fragile]{Multiple seasonal periods}
\protect\phantomsection\label{multiple-seasonal-periods-3}
\begin{Shaded}
\begin{Highlighting}[]
\NormalTok{vic\_elec }\SpecialCharTok{|\textgreater{}} \FunctionTok{gg\_season}\NormalTok{(Demand, }\AttributeTok{period =} \StringTok{"day"}\NormalTok{)}
\end{Highlighting}
\end{Shaded}

\pandocbounded{\includegraphics[keepaspectratio]{01_graphics_files/figure-beamer/unnamed-chunk-7-1.png}}
\end{frame}

\begin{frame}[fragile]{Seasonal subseries plots}
\protect\phantomsection\label{seasonal-subseries-plots}
\begin{itemize}
\tightlist
\item
  Data for each season collected together in time plot as separate time
  series.
\item
  Enables the underlying seasonal pattern to be seen clearly, and
  changes in seasonality over time to be visualised.
\item
  In R: \texttt{gg\_subseries()} (ggtime)
\end{itemize}
\end{frame}

\begin{frame}[fragile]{Quarterly Australian Beer Production}
\protect\phantomsection\label{quarterly-australian-beer-production-2}
\begin{Shaded}
\begin{Highlighting}[]
\NormalTok{beer }\SpecialCharTok{|\textgreater{}} \FunctionTok{gg\_subseries}\NormalTok{(Beer)}
\end{Highlighting}
\end{Shaded}

\pandocbounded{\includegraphics[keepaspectratio]{01_graphics_files/figure-beamer/unnamed-chunk-8-1.pdf}}
\end{frame}

\begin{frame}[fragile]{Australian holidays}
\protect\phantomsection\label{australian-holidays}
\fontsize{9}{10}\sf

\begin{Shaded}
\begin{Highlighting}[]
\NormalTok{holidays }\OtherTok{\textless{}{-}}\NormalTok{ tourism }\SpecialCharTok{|\textgreater{}}
  \FunctionTok{filter}\NormalTok{(Purpose }\SpecialCharTok{==} \StringTok{"Holiday"}\NormalTok{) }\SpecialCharTok{|\textgreater{}}
  \FunctionTok{group\_by}\NormalTok{(State) }\SpecialCharTok{|\textgreater{}}
  \FunctionTok{summarise}\NormalTok{(}\AttributeTok{Trips =} \FunctionTok{sum}\NormalTok{(Trips))}
\end{Highlighting}
\end{Shaded}

\begin{verbatim}
# A tsibble: 640 x 3 [1Q]
# Key:       State [8]
   State Quarter Trips
   <chr>   <qtr> <dbl>
 1 ACT   1998 Q1  196.
 2 ACT   1998 Q2  127.
 3 ACT   1998 Q3  111.
 4 ACT   1998 Q4  170.
 5 ACT   1999 Q1  108.
 6 ACT   1999 Q2  125.
 7 ACT   1999 Q3  178.
 8 ACT   1999 Q4  218.
 9 ACT   2000 Q1  158.
10 ACT   2000 Q2  155.
# i 630 more rows
\end{verbatim}
\end{frame}

\begin{frame}[fragile]{Australian holidays}
\protect\phantomsection\label{australian-holidays-1}
\fontsize{9}{10}\sf

\begin{Shaded}
\begin{Highlighting}[]
\NormalTok{holidays }\SpecialCharTok{|\textgreater{}} \FunctionTok{autoplot}\NormalTok{(Trips) }\SpecialCharTok{+}
  \FunctionTok{labs}\NormalTok{(}\AttributeTok{y =} \StringTok{"thousands of trips"}\NormalTok{, }\AttributeTok{title =} \StringTok{"Australian domestic holiday nights"}\NormalTok{)}
\end{Highlighting}
\end{Shaded}

\pandocbounded{\includegraphics[keepaspectratio]{01_graphics_files/figure-beamer/holidays-plot-1.pdf}}
\end{frame}

\begin{frame}[fragile]{Seasonal plots}
\protect\phantomsection\label{seasonal-plots-2}
\fontsize{9}{10}\sf

\begin{Shaded}
\begin{Highlighting}[]
\NormalTok{holidays }\SpecialCharTok{|\textgreater{}} \FunctionTok{gg\_season}\NormalTok{(Trips) }\SpecialCharTok{+}
  \FunctionTok{labs}\NormalTok{(}\AttributeTok{y =} \StringTok{"thousands of trips"}\NormalTok{, }\AttributeTok{title =} \StringTok{"Australian domestic holiday nights"}\NormalTok{)}
\end{Highlighting}
\end{Shaded}

\includegraphics[width=0.42\linewidth,height=\textheight,keepaspectratio]{01_graphics_files/figure-beamer/graphics1-1.pdf}
\end{frame}

\begin{frame}[fragile]{Seasonal subseries plots}
\protect\phantomsection\label{seasonal-subseries-plots-1}
\fontsize{9}{10}\sf

\begin{Shaded}
\begin{Highlighting}[]
\NormalTok{holidays }\SpecialCharTok{|\textgreater{}} \FunctionTok{gg\_subseries}\NormalTok{(Trips) }\SpecialCharTok{+}
  \FunctionTok{labs}\NormalTok{(}\AttributeTok{y =} \StringTok{"thousands of trips"}\NormalTok{, }\AttributeTok{title =} \StringTok{"Australian domestic holiday nights"}\NormalTok{)}
\end{Highlighting}
\end{Shaded}

\includegraphics[width=\linewidth,height=0.73\textheight,keepaspectratio]{01_graphics_files/figure-beamer/graphics2-1.pdf}
\end{frame}

\begin{frame}[fragile]{Calendar plots}
\protect\phantomsection\label{calendar-plots}
\begin{Shaded}
\begin{Highlighting}[]
\FunctionTok{library}\NormalTok{(ggtime)}
\FunctionTok{library}\NormalTok{(mixtime)}
\NormalTok{vic\_elec }\SpecialCharTok{|\textgreater{}}
  \FunctionTok{filter}\NormalTok{(}\FunctionTok{year}\NormalTok{(Date) }\SpecialCharTok{==} \FunctionTok{year}\NormalTok{(}\DecValTok{2014}\NormalTok{L)) }\SpecialCharTok{|\textgreater{}}
  \FunctionTok{as\_tibble}\NormalTok{() }\SpecialCharTok{|\textgreater{}}
  \FunctionTok{mutate}\NormalTok{(}\AttributeTok{Time =} \FunctionTok{datetime}\NormalTok{(Time)) }\SpecialCharTok{|\textgreater{}}
  \FunctionTok{ggplot}\NormalTok{(}\FunctionTok{aes}\NormalTok{(}\AttributeTok{x =}\NormalTok{ Time, }\AttributeTok{y =}\NormalTok{ Demand)) }\SpecialCharTok{+}
  \FunctionTok{geom\_line}\NormalTok{() }\SpecialCharTok{+}
  \FunctionTok{coord\_calendar}\NormalTok{(}\AttributeTok{rows =} \FunctionTok{month}\NormalTok{(}\DecValTok{1}\NormalTok{L), }\AttributeTok{cols =} \ConstantTok{NULL}\NormalTok{) }\SpecialCharTok{+}
  \FunctionTok{scale\_x\_mixtime}\NormalTok{(}
    \AttributeTok{time\_breaks =} \FunctionTok{days}\NormalTok{(}\DecValTok{1}\NormalTok{L),}
    \AttributeTok{time\_labels =} \StringTok{"\{cyc(day, cal\_isoweek$week, label = TRUE, abbreviate = TRUE)\}"}
\NormalTok{  ) }\SpecialCharTok{+}
  \FunctionTok{theme}\NormalTok{(}\AttributeTok{axis.text.y =} \FunctionTok{element\_blank}\NormalTok{(), }\AttributeTok{axis.ticks.y =} \FunctionTok{element\_blank}\NormalTok{())}
\end{Highlighting}
\end{Shaded}

\begin{itemize}
\tightlist
\item
  \texttt{coord\_calendar()} cuts a time axis into a calendar-like grid
  of rows/columns
\item
  each row above is a month, and each cell within it a day
\end{itemize}
\end{frame}

\begin{frame}{Calendar plots}
\protect\phantomsection\label{calendar-plots-1}
\fontsize{10}{11}\sf

\includegraphics[width=\linewidth,height=0.9\textheight,keepaspectratio]{01_graphics_files/figure-beamer/ggtime_calendar2-1.pdf}
\end{frame}

\section{Seasonal or cyclic?}\label{seasonal-or-cyclic}

\begin{frame}{Time series patterns}
\protect\phantomsection\label{time-series-patterns-2}
\begin{description}
\tightlist
\item[Trend]
pattern exists when there is a long-term increase or decrease in the
data.
\item[Seasonal]
pattern exists when a series is influenced by seasonal factors (e.g.,
the quarter of the year, the month, or day of the week).
\item[Cyclic]
pattern exists when data exhibit rises and falls that are
\emph{not of fixed period} (duration usually of at least 2 years).
\end{description}
\end{frame}

\begin{frame}{Time series components}
\protect\phantomsection\label{time-series-components}
\begin{block}{Differences between seasonal and cyclic patterns:}
\protect\phantomsection\label{differences-between-seasonal-and-cyclic-patterns}
\begin{itemize}
\tightlist
\item
  seasonal pattern constant length; cyclic pattern variable length
\item
  average length of cycle longer than length of seasonal pattern
\item
  magnitude of cycle more variable than magnitude of seasonal pattern
\end{itemize}
\end{block}
\end{frame}

\begin{frame}[fragile]{Time series patterns}
\protect\phantomsection\label{time-series-patterns-3}
\fontsize{10}{10}\sf

\begin{Shaded}
\begin{Highlighting}[]
\NormalTok{aus\_production }\SpecialCharTok{|\textgreater{}}
  \FunctionTok{filter}\NormalTok{(}\FunctionTok{year}\NormalTok{(Quarter) }\SpecialCharTok{\textgreater{}=} \FunctionTok{year}\NormalTok{(}\DecValTok{1980}\NormalTok{L)) }\SpecialCharTok{|\textgreater{}}
  \FunctionTok{autoplot}\NormalTok{(Electricity) }\SpecialCharTok{+}
  \FunctionTok{labs}\NormalTok{(}\AttributeTok{y =} \StringTok{"GWh"}\NormalTok{, }\AttributeTok{title =} \StringTok{"Australian electricity production"}\NormalTok{)}
\end{Highlighting}
\end{Shaded}

\pandocbounded{\includegraphics[keepaspectratio]{01_graphics_files/figure-beamer/unnamed-chunk-10-1.pdf}}
\end{frame}

\begin{frame}[fragile]{Time series patterns}
\protect\phantomsection\label{time-series-patterns-4}
\fontsize{10}{10}\sf

\begin{Shaded}
\begin{Highlighting}[]
\NormalTok{aus\_production }\SpecialCharTok{|\textgreater{}}
  \FunctionTok{autoplot}\NormalTok{(Bricks) }\SpecialCharTok{+}
  \FunctionTok{labs}\NormalTok{(}\AttributeTok{title =} \StringTok{"Australian clay brick production"}\NormalTok{,}
       \AttributeTok{x =} \StringTok{"Year"}\NormalTok{, }\AttributeTok{y =} \StringTok{"million units"}\NormalTok{)}
\end{Highlighting}
\end{Shaded}

\pandocbounded{\includegraphics[keepaspectratio]{01_graphics_files/figure-beamer/unnamed-chunk-11-1.pdf}}
\end{frame}

\begin{frame}[fragile]{Time series patterns}
\protect\phantomsection\label{time-series-patterns-5}
\fontsize{10}{10}\sf

\begin{Shaded}
\begin{Highlighting}[]
\NormalTok{us\_employment }\SpecialCharTok{|\textgreater{}}
  \FunctionTok{filter}\NormalTok{(Title }\SpecialCharTok{==} \StringTok{"Retail Trade"}\NormalTok{, }\FunctionTok{year}\NormalTok{(Month) }\SpecialCharTok{\textgreater{}=} \FunctionTok{year}\NormalTok{(}\DecValTok{1980}\NormalTok{L)) }\SpecialCharTok{|\textgreater{}}
  \FunctionTok{autoplot}\NormalTok{(Employed }\SpecialCharTok{/} \FloatTok{1e3}\NormalTok{) }\SpecialCharTok{+}
  \FunctionTok{labs}\NormalTok{(}\AttributeTok{title =} \StringTok{"Retail employment, USA"}\NormalTok{, }\AttributeTok{y =} \StringTok{"Million people"}\NormalTok{)}
\end{Highlighting}
\end{Shaded}

\pandocbounded{\includegraphics[keepaspectratio]{01_graphics_files/figure-beamer/unnamed-chunk-12-1.pdf}}
\end{frame}

\begin{frame}[fragile]{Time series patterns}
\protect\phantomsection\label{time-series-patterns-6}
\fontsize{10}{10}\sf

\begin{Shaded}
\begin{Highlighting}[]
\NormalTok{gafa\_stock }\SpecialCharTok{|\textgreater{}}
  \FunctionTok{filter}\NormalTok{(Symbol }\SpecialCharTok{==} \StringTok{"AMZN"}\NormalTok{, }\FunctionTok{year}\NormalTok{(Date) }\SpecialCharTok{\textgreater{}=} \FunctionTok{year}\NormalTok{(}\DecValTok{2018}\NormalTok{L)) }\SpecialCharTok{|\textgreater{}}
  \FunctionTok{autoplot}\NormalTok{(Close) }\SpecialCharTok{+}
  \FunctionTok{labs}\NormalTok{(}\AttributeTok{title =} \StringTok{"Amazon closing stock price"}\NormalTok{, }\AttributeTok{x =} \StringTok{"Day"}\NormalTok{, }\AttributeTok{y =} \StringTok{"$"}\NormalTok{)}
\end{Highlighting}
\end{Shaded}

\pandocbounded{\includegraphics[keepaspectratio]{01_graphics_files/figure-beamer/unnamed-chunk-13-1.pdf}}
\end{frame}

\begin{frame}[fragile]{Time series patterns}
\protect\phantomsection\label{time-series-patterns-7}
\fontsize{10}{10}\sf

\begin{Shaded}
\begin{Highlighting}[]
\NormalTok{pelt }\SpecialCharTok{|\textgreater{}}
  \FunctionTok{autoplot}\NormalTok{(Lynx) }\SpecialCharTok{+}
  \FunctionTok{labs}\NormalTok{(}\AttributeTok{title =} \StringTok{"Annual Canadian Lynx Trappings"}\NormalTok{,}
       \AttributeTok{x =} \StringTok{"Year"}\NormalTok{, }\AttributeTok{y =} \StringTok{"Number trapped"}\NormalTok{)}
\end{Highlighting}
\end{Shaded}

\pandocbounded{\includegraphics[keepaspectratio]{01_graphics_files/figure-beamer/unnamed-chunk-14-1.pdf}}
\end{frame}

\begin{frame}{Seasonal or cyclic?}
\protect\phantomsection\label{seasonal-or-cyclic-1}
\alert{Differences between seasonal and cyclic patterns:}

\begin{itemize}
\tightlist
\item
  seasonal pattern constant length; cyclic pattern variable length
\item
  average length of cycle longer than length of seasonal pattern
\item
  magnitude of cycle more variable than magnitude of seasonal pattern
\end{itemize}

\pause

\begin{alertblock}{}
The timing of peaks and troughs is predictable with seasonal data, but unpredictable in the long term with cyclic data.
\end{alertblock}
\end{frame}

\section{Lag plots and
autocorrelation}\label{lag-plots-and-autocorrelation}

\begin{frame}[fragile]{Example: Beer production}
\protect\phantomsection\label{example-beer-production}
\fontsize{10}{11}\sf

\begin{Shaded}
\begin{Highlighting}[]
\NormalTok{new\_production }\OtherTok{\textless{}{-}}\NormalTok{ aus\_production }\SpecialCharTok{|\textgreater{}}
  \FunctionTok{filter}\NormalTok{(}\FunctionTok{year}\NormalTok{(Quarter) }\SpecialCharTok{\textgreater{}=} \FunctionTok{year}\NormalTok{(}\DecValTok{1992}\NormalTok{L))}
\NormalTok{new\_production}
\end{Highlighting}
\end{Shaded}

\begin{verbatim}
# A tsibble: 74 x 7 [1Q]
   Quarter  Beer Tobacco Bricks Cement Electricity   Gas
     <qtr> <dbl>   <dbl>  <dbl>  <dbl>       <dbl> <dbl>
 1 1992 Q1   443    5777    383   1289       38332   117
 2 1992 Q2   410    5853    404   1501       39774   151
 3 1992 Q3   420    6416    446   1539       42246   175
 4 1992 Q4   532    5825    420   1568       38498   129
 5 1993 Q1   433    5724    394   1450       39460   116
 6 1993 Q2   421    6036    462   1668       41356   149
 7 1993 Q3   410    6570    475   1648       42949   163
 8 1993 Q4   512    5675    443   1863       40974   138
 9 1994 Q1   449    5311    421   1468       40162   127
10 1994 Q2   381    5717    475   1755       41199   159
# i 64 more rows
\end{verbatim}
\end{frame}

\begin{frame}[fragile]{Example: Beer production}
\protect\phantomsection\label{example-beer-production-1}
\fontsize{13}{15}\sf

\begin{Shaded}
\begin{Highlighting}[]
\NormalTok{new\_production }\SpecialCharTok{|\textgreater{}} \FunctionTok{gg\_lag}\NormalTok{(Beer)}
\end{Highlighting}
\end{Shaded}

\includegraphics[width=7cm,height=\textheight,keepaspectratio]{01_graphics_files/figure-beamer/unnamed-chunk-16-1.pdf}
\end{frame}

\begin{frame}[fragile]{Example: Beer production}
\protect\phantomsection\label{example-beer-production-2}
\fontsize{13}{15}\sf

\begin{Shaded}
\begin{Highlighting}[]
\NormalTok{new\_production }\SpecialCharTok{|\textgreater{}} \FunctionTok{gg\_lag}\NormalTok{(Beer, }\AttributeTok{geom =} \StringTok{"point"}\NormalTok{)}
\end{Highlighting}
\end{Shaded}

\includegraphics[width=7cm,height=\textheight,keepaspectratio]{01_graphics_files/figure-beamer/unnamed-chunk-17-1.pdf}
\end{frame}

\begin{frame}{Lagged scatterplots}
\protect\phantomsection\label{lagged-scatterplots}
\begin{itemize}
\tightlist
\item
  Each graph shows \(y_t\) plotted against \(y_{t-k}\) for different
  values of \(k\).
\item
  The autocorrelations are the correlations associated with these
  scatterplots.
\item
  ACF (autocorrelation function):

  \begin{itemize}
  \tightlist
  \item
    \(r_1=\text{Correlation}(y_{t}, y_{t-1})\)
  \item
    \(r_2=\text{Correlation}(y_{t}, y_{t-2})\)
  \item
    \(r_3=\text{Correlation}(y_{t}, y_{t-3})\)
  \item
    etc.
  \end{itemize}
\item
  If there is \textbf{seasonality}, the ACF at the seasonal lag (e.g.,
  12 for monthly data) will be \textbf{large and positive}.
\end{itemize}
\end{frame}

\begin{frame}[fragile]{Autocorrelation}
\protect\phantomsection\label{autocorrelation}
Results for first 9 lags for beer data: \fontsize{10}{11}\sf

\begin{Shaded}
\begin{Highlighting}[]
\NormalTok{new\_production }\SpecialCharTok{|\textgreater{}} \FunctionTok{ACF}\NormalTok{(Beer, }\AttributeTok{lag\_max =} \DecValTok{9}\NormalTok{)}
\end{Highlighting}
\end{Shaded}

\begin{verbatim}
# A tsibble: 9 x 2 [1Q]
       lag     acf
  <cf_lag>   <dbl>
1       1Q -0.102 
2       2Q -0.657 
3       3Q -0.0603
4       4Q  0.869 
5       5Q -0.0892
6       6Q -0.635 
7       7Q -0.0542
8       8Q  0.832 
9       9Q -0.108 
\end{verbatim}
\end{frame}

\begin{frame}[fragile]{Autocorrelation}
\protect\phantomsection\label{autocorrelation-1}
Results for first 9 lags for beer data: \fontsize{10}{11}\sf

\begin{Shaded}
\begin{Highlighting}[]
\NormalTok{new\_production }\SpecialCharTok{|\textgreater{}}
  \FunctionTok{ACF}\NormalTok{(Beer, }\AttributeTok{lag\_max =} \DecValTok{9}\NormalTok{) }\SpecialCharTok{|\textgreater{}}
  \FunctionTok{autoplot}\NormalTok{()}
\end{Highlighting}
\end{Shaded}

\pandocbounded{\includegraphics[keepaspectratio]{01_graphics_files/figure-beamer/beeracf-1.pdf}}
\end{frame}

\begin{frame}[fragile]{ACF}
\protect\phantomsection\label{acf}
\begin{Shaded}
\begin{Highlighting}[]
\NormalTok{new\_production }\SpecialCharTok{|\textgreater{}}
  \FunctionTok{ACF}\NormalTok{(Beer) }\SpecialCharTok{|\textgreater{}}
  \FunctionTok{autoplot}\NormalTok{()}
\end{Highlighting}
\end{Shaded}

\pandocbounded{\includegraphics[keepaspectratio]{01_graphics_files/figure-beamer/unnamed-chunk-19-1.pdf}}
\end{frame}

\begin{frame}[fragile]{Australian holidays}
\protect\phantomsection\label{australian-holidays-2}
\fontsize{10}{11}\sf

\begin{Shaded}
\begin{Highlighting}[]
\NormalTok{holidays }\SpecialCharTok{|\textgreater{}} \FunctionTok{ACF}\NormalTok{(Trips)}
\end{Highlighting}
\end{Shaded}

\begin{verbatim}
# A tsibble: 152 x 3 [1Q]
# Key:       State [8]
   State      lag      acf
   <chr> <cf_lag>    <dbl>
 1 ACT         1Q  0.0877 
 2 ACT         2Q  0.252  
 3 ACT         3Q -0.0496 
 4 ACT         4Q  0.300  
 5 ACT         5Q -0.0741 
 6 ACT         6Q  0.269  
 7 ACT         7Q -0.00504
 8 ACT         8Q  0.236  
 9 ACT         9Q -0.0953 
10 ACT        10Q  0.0750 
# i 142 more rows
\end{verbatim}
\end{frame}

\begin{frame}[fragile]{Australian holidays}
\protect\phantomsection\label{australian-holidays-3}
\begin{Shaded}
\begin{Highlighting}[]
\NormalTok{holidays }\SpecialCharTok{|\textgreater{}} \FunctionTok{ACF}\NormalTok{(Trips) }\SpecialCharTok{|\textgreater{}} \FunctionTok{autoplot}\NormalTok{()}
\end{Highlighting}
\end{Shaded}

\includegraphics[width=0.49\linewidth,height=\textheight,keepaspectratio]{01_graphics_files/figure-beamer/tourismacf2-1.pdf}
\end{frame}

\begin{frame}{Trend and seasonality in ACF plots}
\protect\phantomsection\label{trend-and-seasonality-in-acf-plots}
\begin{itemize}
\tightlist
\item
  When data have a trend, the autocorrelations for small lags tend to be
  large and positive.
\item
  When data are seasonal, the autocorrelations will be larger at the
  seasonal lags (i.e., at multiples of the seasonal frequency)
\item
  When data are trended and seasonal, you see a combination of these
  effects.
\end{itemize}
\end{frame}

\begin{frame}[fragile]{US retail trade employment}
\protect\phantomsection\label{us-retail-trade-employment}
\begin{Shaded}
\begin{Highlighting}[]
\NormalTok{retail }\OtherTok{\textless{}{-}}\NormalTok{ us\_employment }\SpecialCharTok{|\textgreater{}}
  \FunctionTok{filter}\NormalTok{(Title }\SpecialCharTok{==} \StringTok{"Retail Trade"}\NormalTok{, }\FunctionTok{year}\NormalTok{(Month) }\SpecialCharTok{\textgreater{}=} \FunctionTok{year}\NormalTok{(}\DecValTok{1980}\NormalTok{L))}
\NormalTok{retail }\SpecialCharTok{|\textgreater{}} \FunctionTok{autoplot}\NormalTok{(Employed)}
\end{Highlighting}
\end{Shaded}

\pandocbounded{\includegraphics[keepaspectratio]{01_graphics_files/figure-beamer/unnamed-chunk-20-1.pdf}}
\end{frame}

\begin{frame}[fragile]{US retail trade employment}
\protect\phantomsection\label{us-retail-trade-employment-1}
\begin{Shaded}
\begin{Highlighting}[]
\NormalTok{retail }\SpecialCharTok{|\textgreater{}}
  \FunctionTok{ACF}\NormalTok{(Employed, }\AttributeTok{lag\_max =} \DecValTok{48}\NormalTok{) }\SpecialCharTok{|\textgreater{}}
  \FunctionTok{autoplot}\NormalTok{()}
\end{Highlighting}
\end{Shaded}

\pandocbounded{\includegraphics[keepaspectratio]{01_graphics_files/figure-beamer/unnamed-chunk-21-1.pdf}}
\end{frame}

\begin{frame}[fragile]{Google stock price}
\protect\phantomsection\label{google-stock-price}
\fontsize{10}{11}\sf

\begin{Shaded}
\begin{Highlighting}[]
\NormalTok{google\_2015 }\OtherTok{\textless{}{-}}\NormalTok{ gafa\_stock }\SpecialCharTok{|\textgreater{}}
  \FunctionTok{filter}\NormalTok{(Symbol }\SpecialCharTok{==} \StringTok{"GOOG"}\NormalTok{, }\FunctionTok{year}\NormalTok{(Date) }\SpecialCharTok{==} \FunctionTok{year}\NormalTok{(}\DecValTok{2015}\NormalTok{L)) }\SpecialCharTok{|\textgreater{}}
  \FunctionTok{select}\NormalTok{(Date, Close)}
\NormalTok{google\_2015}
\end{Highlighting}
\end{Shaded}

\begin{verbatim}
# A tsibble: 252 x 2 [!]
   Date       Close
   <date>     <dbl>
 1 2015-01-02  522.
 2 2015-01-05  511.
 3 2015-01-06  499.
 4 2015-01-07  498.
 5 2015-01-08  500.
 6 2015-01-09  493.
 7 2015-01-12  490.
 8 2015-01-13  493.
 9 2015-01-14  498.
10 2015-01-15  499.
# i 242 more rows
\end{verbatim}
\end{frame}

\begin{frame}[fragile]{Google stock price}
\protect\phantomsection\label{google-stock-price-1}
\begin{Shaded}
\begin{Highlighting}[]
\NormalTok{google\_2015 }\SpecialCharTok{|\textgreater{}} \FunctionTok{autoplot}\NormalTok{(Close)}
\end{Highlighting}
\end{Shaded}

\pandocbounded{\includegraphics[keepaspectratio]{01_graphics_files/figure-beamer/unnamed-chunk-22-1.pdf}}
\end{frame}

\begin{frame}[fragile]{Google stock price}
\protect\phantomsection\label{google-stock-price-2}
\begin{Shaded}
\begin{Highlighting}[]
\NormalTok{google\_2015 }\SpecialCharTok{|\textgreater{}}
  \FunctionTok{ACF}\NormalTok{(Close, }\AttributeTok{lag\_max =} \DecValTok{100}\NormalTok{) }\SpecialCharTok{|\textgreater{}}
  \FunctionTok{autoplot}\NormalTok{()}
\end{Highlighting}
\end{Shaded}

\pandocbounded{\includegraphics[keepaspectratio]{01_graphics_files/figure-beamer/unnamed-chunk-23-1.pdf}}
\end{frame}

\begin{frame}{Which is which?}
\protect\phantomsection\label{which-is-which}
\includegraphics[width=1\linewidth,height=\textheight,keepaspectratio]{01_graphics_files/figure-beamer/unnamed-chunk-24-1.pdf}
\end{frame}

\section{White noise}\label{white-noise}

\begin{frame}[fragile]{Example: White noise}
\protect\phantomsection\label{example-white-noise}
\fontsize{11}{13}\sf

\begin{Shaded}
\begin{Highlighting}[]
\NormalTok{wn }\OtherTok{\textless{}{-}} \FunctionTok{tsibble}\NormalTok{(}\AttributeTok{t =} \FunctionTok{seq}\NormalTok{(}\DecValTok{36}\NormalTok{), }\AttributeTok{y =} \FunctionTok{rnorm}\NormalTok{(}\DecValTok{36}\NormalTok{), }\AttributeTok{index =}\NormalTok{ t)}
\NormalTok{wn }\SpecialCharTok{|\textgreater{}} \FunctionTok{autoplot}\NormalTok{(y)}
\end{Highlighting}
\end{Shaded}

\pandocbounded{\includegraphics[keepaspectratio]{01_graphics_files/figure-beamer/unnamed-chunk-25-1.pdf}}

\only<2>{
\begin{textblock}{10}(1.4,6.7)\fontsize{13}{15}\sf
\begin{alertblock}{}
White noise data is uncorrelated across time with zero mean and constant variance.

(Technically, we require independence as well.)
\end{alertblock}
\end{textblock}}
\end{frame}

\begin{frame}[fragile]{Example: White noise}
\protect\phantomsection\label{example-white-noise-1}
\fontsize{10}{10}\sf

\begin{Shaded}
\begin{Highlighting}[]
\NormalTok{wn }\SpecialCharTok{|\textgreater{}} \FunctionTok{ACF}\NormalTok{(y)}
\end{Highlighting}
\end{Shaded}

\fontsize{10}{10}\sf\tabcolsep=0.1cm

{\def\LTcaptype{none} % do not increment counter
\begin{longtable}[]{@{}
  >{\centering\arraybackslash}p{(\linewidth - 18\tabcolsep) * \real{0.0989}}
  >{\centering\arraybackslash}p{(\linewidth - 18\tabcolsep) * \real{0.0989}}
  >{\centering\arraybackslash}p{(\linewidth - 18\tabcolsep) * \real{0.0989}}
  >{\centering\arraybackslash}p{(\linewidth - 18\tabcolsep) * \real{0.0989}}
  >{\centering\arraybackslash}p{(\linewidth - 18\tabcolsep) * \real{0.0989}}
  >{\centering\arraybackslash}p{(\linewidth - 18\tabcolsep) * \real{0.0989}}
  >{\centering\arraybackslash}p{(\linewidth - 18\tabcolsep) * \real{0.0989}}
  >{\centering\arraybackslash}p{(\linewidth - 18\tabcolsep) * \real{0.0989}}
  >{\centering\arraybackslash}p{(\linewidth - 18\tabcolsep) * \real{0.0989}}
  >{\centering\arraybackslash}p{(\linewidth - 18\tabcolsep) * \real{0.1099}}@{}}
\toprule\noalign{}
\begin{minipage}[b]{\linewidth}\centering
\(r_{1}\)
\end{minipage} & \begin{minipage}[b]{\linewidth}\centering
\(r_{2}\)
\end{minipage} & \begin{minipage}[b]{\linewidth}\centering
\(r_{3}\)
\end{minipage} & \begin{minipage}[b]{\linewidth}\centering
\(r_{4}\)
\end{minipage} & \begin{minipage}[b]{\linewidth}\centering
\(r_{5}\)
\end{minipage} & \begin{minipage}[b]{\linewidth}\centering
\(r_{6}\)
\end{minipage} & \begin{minipage}[b]{\linewidth}\centering
\(r_{7}\)
\end{minipage} & \begin{minipage}[b]{\linewidth}\centering
\(r_{8}\)
\end{minipage} & \begin{minipage}[b]{\linewidth}\centering
\(r_{9}\)
\end{minipage} & \begin{minipage}[b]{\linewidth}\centering
\(r_{10}\)
\end{minipage} \\
\midrule\noalign{}
\endhead
0.018 & 0.120 & -0.171 & -0.232 & 0.025 & 0.080 & -0.148 & 0.215 &
-0.389 & -0.103 \\
\bottomrule\noalign{}
\end{longtable}
}

\pandocbounded{\includegraphics[keepaspectratio]{01_graphics_files/figure-beamer/unnamed-chunk-26-1.pdf}}

\pause\fontsize{13}{15}\sf

\begin{itemize}
\tightlist
\item
  Sample autocorrelations for white noise series.
\item
  Expect each autocorrelation to be close to zero.
\item
  Blue lines show 95\% critical values.
\end{itemize}
\end{frame}

\begin{frame}[fragile]{Example: Pigs slaughtered}
\protect\phantomsection\label{example-pigs-slaughtered}
\fontsize{11}{12}\sf

\begin{Shaded}
\begin{Highlighting}[]
\NormalTok{pigs }\OtherTok{\textless{}{-}}\NormalTok{ aus\_livestock }\SpecialCharTok{|\textgreater{}}
  \FunctionTok{filter}\NormalTok{(State }\SpecialCharTok{==} \StringTok{"Victoria"}\NormalTok{, Animal }\SpecialCharTok{==} \StringTok{"Pigs"}\NormalTok{, }\FunctionTok{year}\NormalTok{(Month) }\SpecialCharTok{\textgreater{}=} \FunctionTok{year}\NormalTok{(}\DecValTok{2014}\NormalTok{L))}
\NormalTok{pigs }\SpecialCharTok{|\textgreater{}} \FunctionTok{autoplot}\NormalTok{(Count }\SpecialCharTok{/} \FloatTok{1e3}\NormalTok{) }\SpecialCharTok{+}
  \FunctionTok{labs}\NormalTok{(}\AttributeTok{x =} \StringTok{"Year"}\NormalTok{, }\AttributeTok{y =} \StringTok{"Thousands"}\NormalTok{,}
       \AttributeTok{title =} \StringTok{"Number of pigs slaughtered in Victoria"}\NormalTok{)}
\end{Highlighting}
\end{Shaded}

\pandocbounded{\includegraphics[keepaspectratio]{01_graphics_files/figure-beamer/unnamed-chunk-27-1.pdf}}
\end{frame}

\begin{frame}[fragile]{Example: Pigs slaughtered}
\protect\phantomsection\label{example-pigs-slaughtered-1}
\begin{Shaded}
\begin{Highlighting}[]
\NormalTok{pigs }\SpecialCharTok{|\textgreater{}}
  \FunctionTok{ACF}\NormalTok{(Count) }\SpecialCharTok{|\textgreater{}}
  \FunctionTok{autoplot}\NormalTok{()}
\end{Highlighting}
\end{Shaded}

\pandocbounded{\includegraphics[keepaspectratio]{01_graphics_files/figure-beamer/unnamed-chunk-28-1.pdf}}
\end{frame}

\begin{frame}{Example: Pigs slaughtered}
\protect\phantomsection\label{example-pigs-slaughtered-2}
Monthly total number of pigs slaughtered in the state of Victoria,
Australia, from January 2014 through December 2018 (Source: Australian
Bureau of Statistics.)\pause

\begin{itemize}
\tightlist
\item
  Difficult to detect pattern in time plot.
\item
  ACF shows significant autocorrelation for lag 2 and 12.
\item
  Indicate some slight seasonality.
\end{itemize}

\pause

These show the series is \textbf{not a white noise series}.
\end{frame}




\end{document}
